%%%%%%%%%%%%%%%%%%%%%%%%%%%%%%%%%%%%%%%%%
% Cheatsheet
% LaTeX Template
% Version 1.0 (12/12/15)
%
% This template has been downloaded from:
% http://www.LaTeXTemplates.com
%
% Original author:
% Michael Müller (https://github.com/cmichi/latex-template-collection) with
% extensive modifications by Vel (vel@LaTeXTemplates.com)
%
% License:
% The MIT License (see included LICENSE file)
%
%%%%%%%%%%%%%%%%%%%%%%%%%%%%%%%%%%%%%%%%%

%----------------------------------------------------------------------------------------
%	PACKAGES AND OTHER DOCUMENT CONFIGURATIONS
%----------------------------------------------------------------------------------------

\documentclass[fontsize=9pt]{scrartcl} % 9pt font size

\usepackage[utf8]{inputenc} % Required for inputting international characters
\usepackage[T1]{fontenc} % Output font encoding for international characters

\usepackage[margin=0pt, landscape]{geometry} % Page margins and orientation

\usepackage{graphicx} % Required for including images

\usepackage{color} % Required for color customization
\usepackage[table]{xcolor}
\definecolor{lightblue}{rgb}{0.145, 0.6666, 1}

\usepackage{url} % Required for the \url command to easily display URLs

\usepackage[ % This block contains information used to annotate the PDF
colorlinks=false, 
pdftitle={IODE v6 Keyboard Shortcuts}, 
pdfauthor={Alix Damman}, 
pdfsubject={IODE v6 Keyboard Shortcuts}, 
pdfkeywords={IODE, keyboard, shortcuts}
]{hyperref}

\setlength{\unitlength}{1mm} % Set the length that numerical units are measured in
\setlength{\parindent}{0pt} % Stop paragraph indentation

\newcommand{\sectiontitle}[1]{\paragraph{#1} \ \\ \rule{\linewidth}{0.2mm} \\} % Custom command for section titles

%----------------------------------------------------------------------------------------

\begin{document}

\rowcolors{1}{white}{lightblue!25} % alternating colors in tables

\begin{picture}(297,210) % Create a container for the page content

%----------------------------------------------------------------------------------------
%	TITLE SECTION 
%----------------------------------------------------------------------------------------

\put(10,200){ % Position on the page to put the title
\begin{minipage}[t]{210mm} % The size and alignment of the title
\section*{\huge IODE v6\newline Keyboard Shortcuts} % Title
\end{minipage}
}

%----------------------------------------------------------------------------------------
%	FIRST COLUMN SPECIFICATION
%----------------------------------------------------------------------------------------

\put(10,180){ % Divide the page
\begin{minipage}[t]{85mm} % Create a box to house text

\sectiontitle{Global Shortcuts}

\begin{tabular}{ p{0.4\textwidth} p{0.6\textwidth} }
    Tab                 & Shift to the next field \\
    Enter               & Start editing field \newline Expand the dropdown menu 
                          \newline Show list of files in the directory 
                          \newline Execute button \\ 
    Esc                 & Stop editing field \newline Exit box dialog \\
    F1                  & Open IODE help \\
    F9                  & Show the used memory \\
    Alt + X             & Exit
\end{tabular}
\newline\newline

\sectiontitle{Menus}

\begin{tabular}{ p{0.4\textwidth} p{0.6\textwidth} }
    $\leftarrow$ / $\rightarrow$  & Shift to the next/previous menu \\
    Enter                         & Expand the menu \newline Open the box dialog \\
    Alt + [a-z]                   & Open the (sub)menu starting with the pressed letter \\
    Esc                           & Close the menu 
\end{tabular}
\newline\newline

\sectiontitle{Box Dialogs}

\begin{tabular}{ p{0.4\textwidth} p{0.6\textwidth} }
    Tab                 & Shift to the next field \\
    Enter               & Start editing field \newline Expand the dropdown menu 
                          \newline Show list of files in the directory 
                          \newline Execute button \\ 
    Esc                 & Stop editing field \newline Exit box dialog \\
    F1                  & Open IODE help \\
    F3                  & Exit \\
    F6                  & Open printing options (if available) \\
    F8                  & Open print box dialog (if available) \\
    F10                 & Press OK 
\end{tabular}
\newline\newline

%----------------------------------------------------------------------------------------

\end{minipage} % End the first column of text
} % End the first division of the page

%----------------------------------------------------------------------------------------
%	SECOND COLUMN SPECIFICATION 
%----------------------------------------------------------------------------------------

\put(105,180){ % Divide the page
\begin{minipage}[t]{85mm} % Create a box to house text

\sectiontitle{Editor Field}

\begin{tabular}{ p{0.4\textwidth} p{0.6\textwidth} }
    Enter                                & Start editing \\  
    Home                                 & Move cursor to the beginning of the line \\
    End                                  & Move cursor to the end of the line \\
    Delete / Backspace                   & Delete the next/previous character \\
    Shift + $\leftarrow$ / $\rightarrow$ & Select next/previous character \\
    Alt + $\leftarrow$ / $\rightarrow$   & Move cursor to the beginning/end of the line \\
    Alt + I                              & Insert line \\
    Alt + J                              & Join next (selected) lines \\ 
    Alt + G                              & Goto line \\ 
    Alt + K                              & Delete current line \\
    Alt + N                              & Split current line \\
    Alt + Q / Esc                        & Stop editing \\
    Alt + M                              & Show all shortcuts \\
    Ctrl + $\leftarrow$ / $\rightarrow$  & Move cursor to the next/previous of the word \\
    Ctrl + Home / End                    & Move cursor to the beginning/end of the editor \\
    Ctrl + Backspace                     & Delete the current word \\ 
    Ctrl + D                             & Delete (cut) block \\
    Ctrl + Y                             & Copy block \\
    Ctrl + C                             & Copy to clipboard \\
    Ctrl + P                             & Paste clipboard \\
    Ctrl + R                             & Find and replace text \\ 
    Ctrl + B                             & select current character \\
    Ctrl + L                             & select current line \\
    Ctrl + U                             & unselect \\
    Ctrl + J                             & Insert line \\
    Ctrl + K                             & Delete current line 
\end{tabular}
\newline\newline

%----------------------------------------------------------------------------------------

\end{minipage} % End the second column of text
} % End the second division of the page

%----------------------------------------------------------------------------------------
%	THIRD COLUMN SPECIFICATION 
%----------------------------------------------------------------------------------------

\put(200,180){ % Divide the page
\begin{minipage}[t]{85mm} % Create a box to house text

\sectiontitle{Workspaces}

\begin{tabular}{ p{0.4\textwidth} p{0.6\textwidth} }
    Space               & Open list of actions \\
    F2                  & Increase size of names column \\
    Shift + F2          & Decrease size of names column \\
    F3                  & Increase size of values column (scalars + variables) \\ 
    Shift + F3          & Decrease size of values column (scalars + variables) \\
    F4                  & Increase the number of displayed decimals (scalars + variables) \\
    Shift + F4          & Decrease the number of displayed decimals (scalars + variables) \\
    F5                  & Change the mode (variables only) \\
    F7                  & (identities) Execute selected identity 
                          \newline (tables) Compute selected table and display values \\
    Shift + F7          & (identities) Open window to execute identities \\
    F8                  & (tables) Compute selected table and plot values as a graph
                          \newline (variables) Display the selected series as a graph \\
    Shift + F8          & (variables) Open window to plot selected variables \\
    Enter               & Start editing IODE object \\
    Insert              & Create a new IODE object \\
    Delete              & Delete the selected IODE object \\
    ~[a-z]              & Go to the first IODE object starting with the pressed letter \\
    Alt + [F1-F7]       & Show the [comments - variables] with the same name or present 
                          in the LEC expression of the current equation/identity \\
    Ctrl + [F1-F7]      & shows the [comments - variables] containing to the current 
                          IODE object \\
    Ctrl + R            & Invert rows and columns (variables only) \\
    Ctrl + Z            & Allow to resize the window of objects \\
    Ctrl + X            & Expand the window of objects to its maximum size \\ 
    Esc                 & Stop displaying workspace 
\end{tabular}
\newline\newline


%----------------------------------------------------------------------------------------

\end{minipage} % End the third column of text
} % End the third division of the page
\end{picture} % End the container for the entire page

%----------------------------------------------------------------------------------------

\begin{picture}(297,210) % Create a container for the page content

%----------------------------------------------------------------------------------------
%	FIRST COLUMN SPECIFICATION
%----------------------------------------------------------------------------------------

\put(10,180){ % Divide the page
\begin{minipage}[t]{85mm} % Create a box to house text

\sectiontitle{Edit Equation + Estimation}

\begin{tabular}{ p{0.4\textwidth} p{0.6\textwidth} }
    F10 & Save equation \\
    Esc & Cancel \\
    F1  & Open IODE help \\
    F3  & Open unit root \\
    F4  & Show estimated coefficients \\
    F5  & Run estimation \\
    F6  & Next equation in block \\
    F7  & Open dynamic adjustment \\
    F8  & Show estimation results menu 
\end{tabular}
\newline\newline

\sectiontitle{Edit Table}

\begin{tabular}{ p{0.4\textwidth} p{0.6\textwidth} }
    F10        & Save table and exit \\
    Esc        & Cancel \\
    F1         & Open IODE help \\
    F2         & Increase size of first column \\
    Shift + F2 & Decrease size of first column \\
    F3         & Increase size of second column \\
    Shift + F3 & Decrease size of second column \\
    F4         & change the attributes of the current cell \\
    F5         & Change the alignment of the current cell (if not LEC) \\
    F6         & Change the font size of the current cell (bold, italic) \\
    F7         & Compute the table and show the values \\
    F8         & Compute and display as a graph \\
    Insert     & Insert a new line \\
    Delete     & Delete a line
\end{tabular}
\newline\newline

%----------------------------------------------------------------------------------------
%	FOOTNOTE
%----------------------------------------------------------------------------------------

\vspace{\baselineskip}
\linethickness{0.5mm} % Thickness of the footer line
{\color{gray}\line(1,0){30}} % Print the line with a custom color

\footnotesize{
Created by Alix Damman, 2026\\ 
\url{https://iode.readthedocs.io/en/stable/index.html}\\
}

%----------------------------------------------------------------------------------------

\end{minipage} % End the first column of text
} % End the first division of the page

\end{picture} % End the container for the entire page

\end{document}
